\section{Static models}
\label{s:models-static}

I use a stylised 3-factor production model throughout this analysis.
$Y$, $A$, $K$, and $L$ are each firm-year observation's GVA, TFP, labour, and capital in levels.
I use lower-case variables to indicate logged terms and bolded variables to indicate $i$-stacked vectors.
\par
Log TFP, denoted by $z_{i,t}$, is the residual after the own-firm contribution of the two classical inputs.
I decompose it to a constant across all observations, firm and time fixed effects and,
most crucially for this investigation, a vector of the weighted sum of all peer 3-factor production factors
with the dot product against the vector of their respective spillover coefficients $\boldsymbol\theta_0$:
\begin{align}
    \mathbf z_t^{spill} &= \begin{pmatrix}
        \sum_{j\neq i} w_{ij}z_{j,t} & \sum_{j\neq i} w_{ij}k_{j,t} & \sum_{j\neq i} w_{ij}l_{j,t}
    \end{pmatrix}
    \\
    \boldsymbol\theta_0 &= \begin{pmatrix}
        \delta_0 & \xi_0 & \zeta_0
    \end{pmatrix}'
\end{align}
%
Different static weights, $w_{ij}$, are explored throughout this section.
The objective is to estimate the vector of structural coefficients $\boldsymbol\theta_0$
of spillover from peers to any given target firm.
\begin{align}
    \label{eq:static-production-3-factor}
    Y_{i,t} &= A_{i,t} \cdot
                K_{i,t}^{\beta_1} \cdot
                L_{i,t}^{\beta_2}
    \\
    \label{eq:static-tfp-identity}
    z_{i,t} &= \ln(A_{it}) = y_{i,t} - \beta_1 k_{i,t} - \beta_2 l_{i,t}
    \\
    \label{eq:static-main-reduced}
            &= \alpha_i + \gamma_t + \beta_0 + \mathbf z_t^{spill}\boldsymbol\theta_0 + \varepsilon_{i,t}
    \\
    \label{eq:static-main-gva}
    y_{i,t} &= \alpha_i + \gamma_t + \beta_0
                + \beta_1 k_{i,t}
                + \beta_2 l_{i,t}
                + \delta_0 \sum_{j\neq i} w_{ij}z_{j,t}
                + \xi_0 \sum_{j\neq i} w_{ij}k_{j,t}
                + \zeta_0 \sum_{j\neq i} w_{ij}l_{j,t}
                + \varepsilon_{i,t}
    \\ \notag
            &
            \begin{aligned}
                =
                \alpha_i + \gamma_t &+ \beta_0 + \beta_1 k_{i,t} + \beta_2 l_{i,t} + \varepsilon_{i,t}
                \\ &
                    + \delta_0 \sum_{j\neq i} w_{ij} \cdot (y_{j,t} - \beta_1 k_{j,t} - \beta_2 l_{j,t})
                    + \xi_0\sum_{j\neq i} w_{ij}k_{j,t}
                    + \zeta_0\cdot \sum_{j\neq i} w_{ij}l_{j,t}                    
            \end{aligned}
    \\
    \label{eq:static-main-peer-effects}
            &
            \begin{aligned}
                =
                \alpha_i + \gamma_t &+ \beta_0 + \beta_1 k_{i,t} + \beta_2 l_{i,t} + \varepsilon_{i,t}
                \\ &
                    + \delta_0 \sum_{j\neq i} w_{ij}y_{j,t}
                    + (\xi_0-\delta_0\beta_1) \cdot \sum_{j\neq i} w_{ij}k_{j,t}
                    + (\zeta_0-\delta_0\beta_2) \cdot \sum_{j\neq i} w_{ij}l_{j,t}                    
            \end{aligned}
\end{align}
\par
Equation \eqref{eq:static-main-peer-effects} is a linear peer effects specification as given by \cite{Manski1993}.
It is helpful to borrow the terminology from this literature think of
$\begin{pmatrix}
    k_{i,t} &   l_{i,t}
\end{pmatrix}\in x_{i,t}$
as \itx{exogenous} characteristics of firm ${i}$ causally affecting its outcome $y_{i,t}$.
There is an \itx{endogenous} effect directly from peer outcomes $y_{j,t}$ with coefficient $\delta_0$.
Peers also have an \itx{exogenous} characteristics $x_{i,t}$,
which have a direct impact on on the target's productivity $x_{i,t}$ via $\xi_0$ and $\zeta_0$.
For this section of models, I assume there are no \itx{correlated} effects\footnote{Though this would be achieved via network fixed effects}.
I use this peer effects terminology throughout this paper italicising \itx{endogenous}, \itx{exogenous}, and \itx{correlated},
noting to the reader the confusion this may cause with the more general and related concepts of endogeneity and exogeneity.
The use of fixed effects has no effect on econometric interpretation of the peer effects model.
\par
It is important to note that I do not impose restrictions on $\boldsymbol\theta_0$.
I allow each of the 3 factors of production to independently exert a spillover effect in general form.
This allows this analysis' investigation as to which of the production factors matters more for productivity spillovers,
and in what ratio, conceptually linked to the discussion in \ref{ss:lit-marshall}
about identifying the channels through which productivity spillovers might occur.
Through this 3-factor production model, it is hard to link each parameter in $\boldsymbol\theta_0$
to a Marshallian channel, but broadly we might consider $\zeta_0$,
the coefficient on an increase in a peer firm's labour inputs affecting a target's TFP to resemble \cite{Marshall1920}'s
labour pooling, with some interaction to knowledge spillovers,
and $\delta_0$, the coefficient from a peer firm's TFP to target TFP to be the remainder of knowledge spillover effects.
A TFP spillover occuring from peer firm's capital $\xi_0$ would be a novel channel to \cite{Marshall1920}'s discussion
and could reflect the idea of more capital-intensive peers .
Regardless, it is strongly justifiable that we would expect the coefficient on each peer production factor
to spillover to a target firm's production heterogenously,
and not be proportionate to the contribution of these factors to their own-firm GVA.
I have no \itx{a priori} arguments for the sign of the $\boldsymbol\theta_0$ coefficients
and their relative magnitudes compared to their contributions to own-firm GVA in $\beta_1, \beta_2$.

\begin{align}
    \label{eq:static-alt-gva-yspill}
    y_{i,t}^{(y)} &= \alpha_i + \gamma_t + \beta_0
                + \beta_1 k_{i,t}
                + \beta_2 l_{i,t}
                + \delta_0 \sum_{j\neq i} w_{ij}y_{j,t}
                + \varepsilon_{i,t}
    \\
    \label{eq:static-alt-gva-zspill}
    y_{i,t}^{(z)} &= \alpha_i + \gamma_t + \beta_0
                    + \beta_1 k_{i,t}
                    + \beta_2 l_{i,t} 
                    + \delta_0 \sum_{j\neq i} w_{ij}z_{j,t}
                    + \varepsilon_{i,t}
    \\ \notag
                &= \alpha_i + \gamma_t + \beta_0
                    + \beta_1 k_{i,t}
                    + \beta_2 l_{i,t} 
                    + \delta_0 \sum_{j\neq i} w_{ij}y_{j,t}
                    - \delta_0\beta_1 \sum_{j\neq i} w_{ij}k_{j,t}
                    - \delta_0\beta_2 \sum_{j\neq i} w_{ij}l_{j,t}
\end{align}
%
Two alternative specifications may be proposed: equation \eqref{eq:static-alt-gva-yspill}
modifies \eqref{eq:static-main-gva} to argue that it is the whole of GVA which matters the peer-target spillover,
not each component heterogenously. It is trivial to see that this implies the coefficient on the \itx{exogenous effect}
is zero, such that $\xi-\delta_0\beta_1=\zeta-\delta_0\beta_2=0$.
This has the interpretation alluded to above,
that the peer effect of any classical factor, capital and labour,
is exactly in the ratio to its own factor income,
which is a strong restriction to impose.
We might expect this to be a nested specification of \eqref{eq:static-main-peer-effects},
that we can merely Wald-test the joint hypothesis as given above,
but actually this specification cannot be identified due to the endogeneity of $y_{j,t}$,
without strict assumptions on the orthogonality of peer errors.
\par
The alternative in equation \eqref{eq:static-alt-gva-zspill} argues that only the TFP component of production
matters for the spillover to influence a target firm's TFP.
This is also a restriction on the more general \eqref{eq:static-main-gva} such that $\xi_0=\zeta_0=0$.
On first glance, it may appear that this model follows the structure of the peer effect class of models too
and is in fact over-identified, with both the coefficients of the peer classical factors and $\sum_{j\neq i}w_{ij}y_{j,t}$
identifying $\delta_0$. However, because of the specific set-up of the reflection problem
where $w_{ij}y_{j,t}$ is endogenous and we instrument it through $k_{j,t},l_{j,t}$ higher-order spatial lags,
this model is actually non-identified. I demonstrate this in \nameref{app:model-zspill-reflection}.
\par
While the classical production function \eqref{eq:static-main-gva} is limited to constant returns to scale,
$\beta_1 + \beta_2\leq 0$, a standard assumption across the growth and firm production literature,
no such restrictions for the firm occur with the endogenous decomposition of TFP.
A restriction on logged production \eqref{eq:static-main-gva} collectively on $\beta_1$, $\beta_2$, and $\boldsymbol\theta_0$
is difficult to justify.
As with other external effects, there is no reason to believe that market factor prices
and the firm decision-rule for spatial allocation in view of these spillover effects, $w_{ij}z_{j,t}$, is socially optimal.
A peer firm $j$ with extremely high and growing components of GVA has little interest in siting themselves
in the centre of a cluster of firms that are struggling with TFP growth,
yet the external impacts of that firm would be significant.

% ------

\subsection{Static panel identification}
\label{ss:static-identification}

Consistency requires the plausible exogeneity of regressors
used in structural equation \eqref{eq:static-main-peer-effects}
from the error term $\varepsilon_{i,t}$.
Capital and labour inputs in every firm $i,j$ are assumed to be exogenous to every
other firm's TFP error shock in this static model.
This is crucial for consistency with the direct own-firm $k_{i,t}, l_{i,t}$
inputs, but also the use of $\mathbf w_{ij} k_t, \mathbf W^k l_t$
as an \itx{exogenous} effect and, as we will later see,
our ability to leverage higher-order spatial lags of the \itx{exogenous} effect
as an unconfounded instrument for $w_{ij}y_{j,t}$.
Broadly, we can interpret this as meaning
factor input decisions are predetermined for this section;
capital faces some time-to-build constraint
where each firm allocates investment in the previous period
based on the limited information available then, $\Omega_{t-1}$.
A similar pattern is assumed of sticky labour decisions,
perhaps through job searching and hiring frictions.
These restrictive assumptions are relaxed in \nameref{s:models-dynamic}.

\begin{align}
    % \notag
    % x_{i,t} &=
    % \begin{pmatrix}
    %     k_{i,t} &   l_{i,t}
    % \end{pmatrix}
    % \\
    \label{eq:static-k-pim}
    K_{i,t} &= I_{i,t} + dK_{i,t-1}
    \\
    \label{eq:static-inv-information}
    I_{i,t} &= \mathbb E[I_{i,t}^*|\Omega_{t-1}]
    \\
    \label{eq:static-l-information}
    L_{i,t} &= \mathbb E[L_{i,t}^*|\Omega_{t-1}]
    \\
    \label{eq:static-predetermined-factors}
    K_{i,t}, L_{i,t} &\in \Omega_{t-1}
    \\ \notag
    \mathbb E[\varepsilon_{i,t}|\Omega_{t-1}] &= 0 
    \\
    \label{eq:static-identification-factor-exogeneity}
    \mathbb E[
        \varepsilon_{i,t}
        \vert
        \begin{pmatrix}
            \mathbf k_t \\ \mathbf l_t
        \end{pmatrix}
    ] &= 0_{(2\times 1)}
    ,\ \forall i
\end{align}

The \itx{endgenous} effect regressor $\sum_{j\neq i} w_{ij} y_{j,t}$ is, unsurprisingly, endogenous,
making valid estimation in this setting more complicated.
To show the endogeneity requires stacking $w_{ij}$ target-peer weights to an $n\times n$ interactions matrix, $\textbf W$.
I invert the \itx{endogenous} component to give a $(\mathbb I-\delta_0\mathbf W)^{-1}$ coefficient on all variables.
\begin{equation}
    \delta_0 < 1 \implies \vert \mathbb I-\delta_0 \mathbf W \vert \neq 0
\end{equation}
The inversion is permissible because I can expect a restriction to apply, ruling out iteratively increasing shocks\footnote{
    $\delta_0\geq1$ would imply an exogenous TFP shock to firm $k$ would
    increase TFP in its peer firm $j$ more than 1x,
    which would increase TFP in that firm's peer $i$ more than 1x,
    and so on until the effect multiplies to infinity,
    or at least on the knife-edge $\delta_0=1$ case,
    any TFP shock occurring in one firm applies equally
    and consistently to every firm in a network.
}. We can interpret this as a `spatial lag polynomial', directly comparable
to the lag polynomial and MA inversion procedure in autoregressive time series analysis\footnote{
    A full algebraic derivation from lines \eqref{eq:static-reflection-problem} to \eqref{eq:static-red-endogenous-effect}
    is available in Appendix \ref{app:model-zspill-reflection}.
    Note that since the vectors $\mathbf k_t, \mathbf l_t$ appear twice in \eqref{eq:static-reflection-problem},
    as components of $\mathbf y_t$ the coefficients can be grouped
    but it is clearer to separate this into an 'own-firm'
    effect of $\mathbf k_t, \mathbf l_t$ and a spillover effect of
    $\mathbf W^{k+1}k_t, \mathbf W^{k+1}l_t\ \forall k\geq 0$
}.
\begin{align}
    \label{eq:static-vector-reduced}
    \mathbf y_t &= \boldsymbol \alpha
                    + \gamma_t \boldsymbol \iota
                    + \beta_0 \boldsymbol \iota
                    + \beta_1 \mathbf k_{t}
                    + \beta_2 \mathbf l_{t}
                    + \delta_0 \mathbf W \mathbf y_t
                    + (\xi_0-\delta_0\beta_1) \mathbf W \mathbf k_t
                    + (\zeta_0-\delta_0\beta_2) \mathbf W \mathbf l_t 
                    + \boldsymbol \varepsilon_{t}
    \\ \notag
            &= (\mathbb I - \delta_0 \mathbf W)^{-1} (
                    \boldsymbol \alpha
                    + \gamma_t \boldsymbol \iota
                    + \beta_0 \boldsymbol \iota
                    + \beta_1 \mathbf k_{t}
                    + \beta_2 \mathbf l_{t}
                    + (\xi_0-\delta_0\beta_1) \mathbf W \mathbf k_t
                    + (\zeta_0-\delta_0\beta_2) \mathbf W \mathbf l_t 
                    + \boldsymbol \varepsilon_{t}
                )
    \\
            &= \sum_{k=0}^\infty \delta_0^k \mathbf W^k (
                    \boldsymbol \alpha
                    + \gamma_t \boldsymbol \iota
                    + \beta_0 \boldsymbol \iota
                    + \beta_1 \mathbf k_{t}
                    + \beta_2 \mathbf l_{t}
                    + (\xi_0-\delta_0\beta_1) \mathbf W \mathbf k_t
                    + (\zeta_0-\delta_0\beta_2) \mathbf W \mathbf l_t 
                    + \boldsymbol \varepsilon_{t}
                )
    \\[1.5ex]
    \label{eq:static-reflection-problem}
            & \begin{aligned}
                =   \sum_{k=0}^\infty \delta_0^k \mathbf W^k
                        \boldsymbol \alpha
                    &+\frac{\gamma_t + \beta_0}{1-\delta_0}\boldsymbol \iota
                    + \beta_1 \mathbf k_{t}
                    + \beta_2 \mathbf l_{t}
                    \\ &
                    + \xi_0\sum_{k=0}^\infty \delta_0^k \mathbf W^{k+1} \mathbf k_t
                    + \zeta_0\sum_{k=0}^\infty \delta_0^k \mathbf W^{k+1} \mathbf l_t                    
                    + \sum_{k=0}^\infty \delta_0^k \mathbf W^k
                        \boldsymbol \varepsilon_{t}
            \end{aligned}
    \\[1.5ex]
    \label{eq:static-red-endogenous-effect}
    y_{j,t}     &= \dots
                    + \underbrace{
                        \sum_{k=0}^\infty
                        \left[
                            \delta_0^k \mathbf W^{k+1}
                        \right]_{(j,)}
                        \begin{pmatrix}\mathbf k_t & \mathbf l_t\end{pmatrix}
                        \begin{pmatrix}
                            \xi_0       \\
                            \zeta_0
                        \end{pmatrix}
                    }_{\text{exogenous}_j}
                    + \underbrace{
                        \sum_{k=0}^\infty
                        [\delta_0^k \mathbf W^k]_{(j,)} \cdot
                    \boldsymbol \varepsilon_t
                    }_{\text{error transformation}_j}
\end{align}

It is the primary econometric challenge of this section to overcome the endogeneity of $y_{j,t}$,
which is clearly demonstrated in reduced-form equation \eqref{eq:static-red-endogenous-effect}.
It contains a linear transformations of all firms' error shocks $\boldsymbol \varepsilon_{t}$,
including $\varepsilon_{i,t}$. This implies:
\begin{equation}
    \mathbb E[y_{j,t}\varepsilon_{i,t}]\neq 0
\end{equation}
\par
The `Reflection Problem' poses a second issue as outlined in \cite{Manski1993}.
Equation \eqref{eq:static-reflection-problem} shows that $\mathbf W \mathbf y_t$,
is a linear combination of all other regressors. In the classical setup,
$\mathbf W$ is a block diagonal matrix of equal-size groups,
and \cite{Manski1993}, later reaffirmed in \cite{Bramoulle2009}
show that the reduced-form estimation of this model contains fewer estimated parameters
than structural parameters, making variation in $y_{j,t}$ indistinguishable
to that of all other regressors.
\par
There are two solutions. I could either proceed with an extremely strong exogeneity assumption,
\eqref{eq:static-ass-ma-exogeneity}, which is most likely satisfied by total orthogonality
between all contemperaneous idiosyncratic errors hitting peer firms \eqref{eq:static-ass-error-orthogonality}.
This in turn makes the regressor $y_{j,t}$ exogenous,
allowing direct estimation of \eqref{eq:static-main-gva}, sidestepping the `Reflection Problem'
by assuming away \itx{endogenous} regressors, and having each regressor exogenously identified in itself.
For identification, this implies that any productivity shock $\varepsilon_{j,t}$ occurring in peer firm $j$,
if it is contemperaneously correlated with TFP ${z_i,t}$ in target firm $i$
must be attributable to the peer effect coefficient via  $\delta_0$.
This precludes any contemperaneous, correlated independent TFP shocks,
which might occur on both the global and local level.
Model (1.) takes this assumption in \ref{ss:model-lmm}.
We should assess these results with the heavy caveat that this strict error orthogonality
assumption is unlikely to be credible.
\begin{align}
    \notag
    \mathbb E[
        \varepsilon_i
        \vert
        \boldsymbol W^k
        \begin{pmatrix}
            \boldsymbol\alpha           &
            \gamma_t \boldsymbol\iota   &
            \beta_0\boldsymbol\iota     &
            \boldsymbol k_t             &
            \boldsymbol l_t             
        \end{pmatrix}
    ] &=0,\  \forall k
    \\
    \label{eq:static-ass-ma-exogeneity}
    \mathbb E[
        \varepsilon_{i,t}
        \vert
        \boldsymbol W^k
        \boldsymbol \varepsilon_t
    ] &=0,\  
    \forall i,j
    \\
    \label{eq:static-ass-error-orthogonality}
    \impliedby \mathbb E[\varepsilon_{i,t}\varepsilon_{j,t}] &= 0,\ \forall k
\end{align}

Instead, as recommended in \cite{Bramoulle2009},
the most defensible approach is to internall instrument the \itx{endogenous} variable, $y_{j,t}$
through deeper spatial lags of the \itx{exogenous} parameters
$\mathbf W^2\mathbf x_t, \mathbf W^3\mathbf x_t \dots \mathbf W^k \mathbf x_t\ \forall k>2$.
These are valid instruments under certain testable properties of $\mathbf W$ given in \tablerefp{tab:bdf_identification}.
Essentially, these reduce to checking that the matrices $\mathbb I, \mathbf W, \mathbf W^2$
are linearly independent, which are sufficiently achieved for group interactions by
having 2 or more differently-sized groups, or a network diameter of at least 2.
\par
This identification strategy is not a `free lunch'.
It requires the underlying assumption that the group size, or network diameter,
is in itself exogenously determined. This would be intuitively violated,
if firms colate based on the expected productivity benefits of moving to an area,
namely to benefit from the supply-side spillovers we aim to investigate.
Unfortunately, it is highly likely that firm location decisions are endogenous,
this is the whole principle of the \nameref{ss:lit-neg} literature whose
descriptive patterns of firm movements underlying productivity spillover I here aim to investigate.
I am helped here, however, by the use of time-invariant weights.
Plausibly, over the observations of a panel dataset up to 20 years,
if firm weights apply to each firm on a fixed basis,
the temporal variation despite these constant location effects provide credible exogenous variation.
On this however, the strategy runs into the limitations of Fame location data discussed in \ref{ss:data-location-limits}
The analysis is also helped by the relatively higher number of panel observations
to relatively low ratio of confounding movers in the dataset,
as well as the numerous, plausibly exogenous reasons why a firm might choose 
to set up in a location with path dependency, rather than to fully maximise
its peer spillover benefits.
On the balance of arguments, I take the assumption of network exogeneity
and identification approach for models in Tables \todo{XII and IX}.
\par
I use the instruments
\todo{$\mathbf W^2\mathbf x_t$ and other second-order variants} for this section's group models \todo{Table VI, models (5-8)}
and $\mathbf W^2\mathbf x_t, \mathbf W^3\mathbf x_t, \mathbf W^4$ as feasible instruments for this section's network models
\todo{Table VII, model}. Equation \eqref{eq:static-moment-conditions} illustrates the valid $2l$ moment conditions,
noting that I have chosen to display this in a $l \times 2$ matrix for space and conceptual symmetry,
as the instruments can be grouped by $k_t$ and $l_t$ regressors.
The first 2 rows correspond to the exogenous regressors within the reduced-form equation itself,
the own-firm classical factors effect from $k_{i,t},l_{i,t}$
and the \itx{exogenous} peer effect from $k_{j,t},l_{j,t}$.
The remaining rows are instrument linear combinations of all
firm classical factor inputs $\mathbf k_t, \mathbf l_t$,
using the higher-order spatial lags of $\mathbf W$ discussed.
\begin{align}
    \label{eq:static-moment-conditions}
    \mathbb E[
        \epsilon_{i,t}\ 
        \vert
        \begin{pmatrix}
            k_{i,t}                                     & l_{i,t} \\
            w_{ij}k_{j,t}                               & w_{ij}l_{j,t}   \\
            \left[\mathbf W^2\right]_{(j,)} \mathbf k_t & \left[\mathbf W^2\right]_{(j,)} \mathbf l_t\\
            \left[\mathbf W^3\right]_{(j,)} \mathbf k_t & \left[\mathbf W^3\right]_{(j,)} \mathbf l_t\\
            \vdots & \vdots \\
            \left[\mathbf W^k\right]_{(j,)} \mathbf k_t & \left[\mathbf W^k\right]_{(j,)} \mathbf l_t
        \end{pmatrix}
    ]
    &= 0_{l\times 2}
\end{align}

These instruments and conditions are intuitive:
we can see through equation \eqref{eq:static-red-endogenous-effect} that
a linear combination of all firms' contemperaneous capital and labour inputs,
$\mathbf W^{k+1}\begin{pmatrix}\mathbf k_t & \mathbf l_t\end{pmatrix}\ \forall k\geq0$ 
are an exogenous component of $y_{j,t}$ with the coefficients $\xi_0$ and $\zeta_0$ respectively.
It has relevance, meaning $\mathbb E[y_{j,t}\cdot [\mathbf W^k]_{(j,)}\mathbf x_t]\neq 0,\ \forall k>2$
as long as $\xi_0\neq 0$ and $\zeta_0 \neq 0$.
The relevance conditions themselves are also unsurprising and highly intuitive
given the model's setup with GVA equation \eqref{eq:static-main-gva}.
Note that the relevance conditions are more complex setup in \cite{Bramoulle2009},
this is because the factor $\delta_0 \begin{pmatrix}\beta_1 & \beta_2\end{pmatrix}$
cancels out due the accounting composition of TFP in the \itx{exogenous effect}, which is absent from the general model.
This also allows clear understanding of why the TFP-only spillover model \eqref{eq:static-alt-gva-zspill} fails,
as setting both $\xi_0=\zeta_0=0$ means the higher order spatial lags of $\mathbf k_t$ and $\mathbf l_t$
disappears from the composition of the $y_{j,t}$ term.
An Anderson-Rubin test for weak instruments in my nested 3-factor spillover setup
can test for the the strength of instrument of $\mathbf k_t$ and $\mathbf l_t$.
If I find that they the higher-order spatial lags are jointly weak instruments,
that would validate the TFP-only spillover model and suggest $\xi_0=\zeta_0=0$.
\par
Finally, one might observe that $[\mathbf W]_{(j,)}\mathbf x_t$
is a component of $y_{j,t}$ and so a valid instrument.
However, it is does not increase the number of instruments I would use in my instrumental specification
as it is already included as an exogenous regressor in \ref{eq:static-vector-reduced}
so does not help increase our identification of $y_{j,t}$.

% --- %

\subsection{Linear-in-means grouped model}
\label{ss:model-lmm}

The first model proceeds with the strict error orthogonality
\eqref{eq:static-ass-ma-exogeneity} and \eqref{eq:static-ass-error-orthogonality}
assumptions. Estimation has the following reduced form:
\begin{align}
    \eqreref{eq:static-main-peer-effects}
    y_{i,t} &= \alpha_i + \gamma_t + \beta_0 + \beta_1 k_{i,t} + \beta_2 l_{i,t} + \varepsilon_{i,t}
        \\ \notag
            &\qquad
                    + \delta_0 \sum_{j\neq i} w_{ij}y_{j,t}
                    + (\xi_0-\delta_0\beta_1) \cdot \sum_{j\neq i} w_{ij}k_{j,t}
                    + (\zeta_0-\delta_0\beta_2) \cdot \sum_{j\neq i} w_{ij}l_{j,t}
    \\ \label{eq:static-lmm-reduced-estimated}
    y_{i,t} &= \alpha_i + \gamma_t + \beta_0 + \beta_1 k_{i,t} + \beta_2 l_{i,t} + \varepsilon_{i,t}
                    + \beta_3 \sum_{j\neq i} w_{ij}y_{i,t}
                    + \beta_4 \sum_{j\neq i} w_{ij}k_{j,t}
                    + \beta_5 \sum_{j\neq i} w_{ij}l_{j,t}
    \\
    \mathbb E[ &\varepsilon_{i,t} \cdot
        \begin{pmatrix}
            k_{i,t} &
            l_{i,t} &
            y_{j,t} &
            k_{j,t} &
            l_{j,t}
        \end{pmatrix}'
    ] = 0_{(5\times 1)}
    \\ &\qquad
    \begin{pmatrix}
        \hat\beta_1     \\
        \hat\beta_2     \\
        \hat\beta_3     \\
        \hat\beta_4     \\
        \hat\beta_5
    \end{pmatrix}
    =
    \begin{pmatrix}
        \beta_1     \\
        \beta_2     \\
        \delta_0     \\
        \xi_0-\delta_0\beta_1     \\
        \zeta_0-\delta_0\beta_2
    \end{pmatrix}
\end{align}

With 5 structural and 5 reduced-form parameters, this model is identified.
The initial static model applies simple time-invariant weights $w_{ij}$
in a group structure, $\mathbf W_{g1}$
As has been indicated by the sum of peers $j\neq i$,
I adopt a 'leave-one-out' approach in line with
the zero-diagonal structure of $\mathbf W_{g1}$ described.
Intuitively, this is the unweighted average of peer TFP within a spatial ring.
This weighting matrix is comparable to many sets of weights applied for the 'Linear-in-Means (LMM)' Model
common to the peer effects literature, for instance \cite{Moffitt2001}.
I discussed group assignment in section \ref{ss:data-distance}.
The initial Model (1.) estimates \eqref{eq:static-lmm-reduced-estimated},
meaning a single group $\mathbf W_{g1}$ corresponding to the \itx{pc8} group,
assumeing the full vector of error orthogonality implying regressor exogeneity.

\begin{align}
    \label{eq:static-lmm-iv-fd-reduced-estimated}
    \Delta y_{i,t} &= \Delta\gamma_t + \beta_1 \Delta k_{i,t} + \beta_2 \Delta l_{i,t} + \Delta \varepsilon_{i,t}
                    + \beta_3 \Delta \sum_{j\neq i} w_{ij}y_{i,t}
                    + \beta_4 \Delta \sum_{j\neq i} w_{ij}k_{j,t}
                    + \beta_5 \Delta \sum_{j\neq i} w_{ij}l_{j,t}
    \\
    \eqreref{eq:static-moment-conditions}
    \mathbb E &[
        \epsilon_{i,t}\ 
        \vert
        \begin{pmatrix}
            k_{i,t}                                     & l_{i,t} \\
            w_{ij}k_{j,t}                               & w_{ij}l_{j,t}   \\
            \left[\mathbf W^2\right]_{(j,)} \mathbf k_t & \left[\mathbf W^2\right]_{(j,)} \mathbf l_t\\
            \left[\mathbf W^3\right]_{(j,)} \mathbf k_t & \left[\mathbf W^3\right]_{(j,)} \mathbf l_t\\
            \vdots & \vdots \\
            \left[\mathbf W^k\right]_{(j,)} \mathbf k_t & \left[\mathbf W^k\right]_{(j,)} \mathbf l_t
        \end{pmatrix}
    ]
    = 0_{l\times 2}
\end{align}

Model (2) instruments these groups via $\mathbf W^2\mathbf x$,
in line with the moment conditions in \eqref{eq:static-moment-conditions}.
As an extension in \ref{tab:lmm-3r}, I regress the multiple, non-overlapping spatial groups at once.
I add $\mathbf W_{g1}, \mathbf W_{g2}, \mathbf W_{g3}$, concentric doughnut groups
for \itx{pc4} and \itx{ttwa} to run all in a panel OLS regression simultaneously.

\begin{table}[hbtp]
    \centering
    \begin{threeparttable}
        \caption{LMM: 1 peer regressor with group weights, error orthogonality model vs instrumented}%
        \label{tab:lmm-pc8}
        	\begin{tabular}{lcc}
		\toprule\toprule
		\lblock{
			~
			\\
			~
		} & 		\cblock{
			\footnotesize{(1.)}
			\\
			\small{\itx{y}}
		} & 		\cblock{
			\footnotesize{(2.)}
			\\
			\small{\itx{y}}
		} \\[0.8em]
		\toprule
		\lblock{
			constant
			\\
			~
		} & 		\cblock{
			$0.006$
			\\
			\footnotesize{(0.001)***}
		} &  
		\\ [0.9em]
		\lblock{
			\itx{k}
			\\
			~
		} & 		\cblock{
			$0.238$
			\\
			\footnotesize{(0.007)***}
		} & 		\cblock{
			$0.233$
			\\
			\footnotesize{(0.008)***}
		}
		\\ [0.9em]
		\lblock{
			\itx{l}
			\\
			~
		} & 		\cblock{
			$0.552$
			\\
			\footnotesize{(0.007)***}
		} & 		\cblock{
			$0.555$
			\\
			\footnotesize{(0.007)***}
		}
		\\ [0.9em]
		\lblock{
			\itx{wg1\_y}
			\\
			~
		} & 		\cblock{
			$0.171$
			\\
			\footnotesize{(0.004)***}
		} & 		\cblock{
			$-0.384$
			\\
			\footnotesize{(0.177)**}
		}
		\\ [0.9em]
		\lblock{
			\itx{wg1}\textsubscript{$k$}
			\\
			~
		} & 		\cblock{
			$-0.050$
			\\
			\footnotesize{(0.003)***}
		} & 		\cblock{
			$0.118$
			\\
			\footnotesize{(0.054)**}
		}
		\\ [0.9em]
		\lblock{
			\itx{wg1\_l}
			\\
			~
		} & 		\cblock{
			$-0.101$
			\\
			\footnotesize{(0.004)***}
		} & 		\cblock{
			$0.249$
			\\
			\footnotesize{(0.112)**}
		}
		\\ [0.9em]
		\hline
		\textit{t} fixed effects & \checkmark & ~
		\\
		Observations & 587,772 & 587,772
		\\
		R$^2$ (excl. F.E.) & 0.1336 & 0.0012
		\\
		R$^2$ (Overall) & 0.1350 & 
		\\
		\bottomrule
	\end{tabular}
    \end{threeparttable}
    \\~\\
    \vspace{1em}
    \begin{threeparttable}
        \caption{LMM: 3 non-overlapping groups, error orthogonality vs instrumented specification}
        \label{tab:lmm-3r}
        	\begin{tabular}{lcc}
		\toprule\toprule
		\lblock{
			~
			\\
			~
		} & 		\cblock{
			\footnotesize{(3.)}
			\\
			\small{\itx{y}}
		} & 		\cblock{
			\footnotesize{(4.)}
			\\
			\small{\itx{y}}
		} \\[0.8em]
		\toprule
		\lblock{
			constant
			\\
			~
		} & 		\cblock{
			$2.342$
			\\
			\footnotesize{(0.095)***}
		} &  
		\\ [0.9em]
		\lblock{
			\itx{k}
			\\
			~
		} & 		\cblock{
			$0.261$
			\\
			\footnotesize{(0.005)***}
		} & 		\cblock{
			$0.236$
			\\
			\footnotesize{(0.011)***}
		}
		\\ [0.9em]
		\lblock{
			\itx{l}
			\\
			~
		} & 		\cblock{
			$0.649$
			\\
			\footnotesize{(0.005)***}
		} & 		\cblock{
			$0.580$
			\\
			\footnotesize{(0.020)***}
		}
		\\ [0.9em]
		\lblock{
			\itx{wg1\_y}
			\\
			~
		} & 		\cblock{
			$0.198$
			\\
			\footnotesize{(0.004)***}
		} & 		\cblock{
			$-0.715$
			\\
			\footnotesize{(0.596)}
		}
		\\ [0.9em]
		\lblock{
			\itx{wg1}\textsubscript{$k$}
			\\
			~
		} & 		\cblock{
			$-0.058$
			\\
			\footnotesize{(0.003)***}
		} & 		\cblock{
			$0.215$
			\\
			\footnotesize{(0.180)}
		}
		\\ [0.9em]
		\lblock{
			\itx{wg1\_l}
			\\
			~
		} & 		\cblock{
			$-0.117$
			\\
			\footnotesize{(0.004)***}
		} & 		\cblock{
			$0.462$
			\\
			\footnotesize{(0.379)}
		}
		\\ [0.9em]
		\lblock{
			\itx{wg2\_y}
			\\
			~
		} & 		\cblock{
			$0.019$
			\\
			\footnotesize{(0.006)***}
		} & 		\cblock{
			$4.312$
			\\
			\footnotesize{(3.073)}
		}
		\\ [0.9em]
		\lblock{
			\itx{wg2}\textsubscript{$k$}
			\\
			~
		} & 		\cblock{
			$-0.012$
			\\
			\footnotesize{(0.004)***}
		} & 		\cblock{
			$-1.493$
			\\
			\footnotesize{(1.088)}
		}
		\\ [0.9em]
		\lblock{
			\itx{wg2\_l}
			\\
			~
		} & 		\cblock{
			$-0.001$
			\\
			\footnotesize{(0.007)}
		} & 		\cblock{
			$-2.833$
			\\
			\footnotesize{(2.002)}
		}
		\\ [0.9em]
		\lblock{
			\itx{wg3\_l}
			\\
			~
		} & 		\cblock{
			$-0.012$
			\\
			\footnotesize{(0.019)}
		} & 		\cblock{
			$26.034$
			\\
			\footnotesize{(16.288)}
		}
		\\ [0.9em]
		\lblock{
			\itx{wg3}\textsubscript{$k$}
			\\
			~
		} & 		\cblock{
			$-0.021$
			\\
			\footnotesize{(0.012)*}
		} & 		\cblock{
			$14.842$
			\\
			\footnotesize{(9.256)}
		}
		\\ [0.9em]
		\lblock{
			\itx{wg3\_y}
			\\
			~
		} & 		\cblock{
			$0.048$
			\\
			\footnotesize{(0.017)***}
		} & 		\cblock{
			$-40.411$
			\\
			\footnotesize{(25.266)}
		}
		\\ [0.9em]
		\hline
		\textit{i} fixed effects & \checkmark & ~
		\\
		\textit{t} fixed effects & \checkmark & ~
		\\
		Observations & 718,743 & 581,185
		\\
		R$^2$ (excl. F.E.) & 0.3543 & -11.504
		\\
		R$^2$ (Overall) & 0.7970 & 
		\\
		\bottomrule
	\end{tabular}
        \begin{tablenotes}
            All parameters reported with standard errors underneath.
            Significant estimates against an $H_0: \beta=0$ are indicated
            with $p<0.1$ (*), $p<0.05$ (**), and $p<0.01$ (***).
        \end{tablenotes}
    \end{threeparttable}
\end{table}

\par
I can show algebraically that multiple spatial autoregressors do not change the reflection problem:
the moving average inversion's coefficient merely becomes the weighted sum
of 3 independent matrices $\theta_g \mathbf W_g$.
As long as there is sufficient heterogeneity between the weighting matrices,
which in these cases is guaranteed by the orthogonality of $\mathbf W_g$
matrices to each other, then there is no risk of multicollinearity.
Coefficients in the results $\boldsymbol \hat\theta_1$ for the closest spatial group \itx{pc8}
are identical to $\boldsymbol \hat\theta_0$ in the one-group model.
This is an unsurprising result justified by the orthogonality, since the additional regressors
do not affect the variation attributable to the $\mathbf W_1$-lagged variables
as the regressors with $\mathbf W_2, \mathbf W_3$ are orthogonal.
\todo{Modify equation 18 to include $\sum W^3$ in the inversion term}


% ---
\newpage
\subsection{Static distance decay}
\label{ss:model-dd}

To improve on the LMM specification, I make a modification where
every peer firm is weighted by its influence on the target by the reciprocal of the 
continuous distance metric calculated in \ref{ss:data-distance}
$w_{ijt}=f(d)$. This is stylistic of a declining peer influence further away,
representing a \todo{lower frequency of social collisions, barriers created via transport and time costs,
and an increased difficulty of observation of management practices}.

\begin{align}
    \label{eq:dd-structural1}
    z_{it} &= \alpha_i + \gamma_t + \beta_0
            + \beta_1 \sum_{j\neq i} w_{ijt}z_{jt} + \varepsilon_{it}
    \\
    \mathbf z_{t} &= (\alpha_i + \gamma_t + \beta_0)\boldsymbol\iota
            + \beta_1 \sum_{j\neq i} \Omega_t \mathbf z_{t} + \boldsymbol\varepsilon_{t}
    \\
    w_{ijt}\in \mathbf \Omega_t &=
        \begin{cases}
            \frac1{d_{ij}},         &\text{(1) gravity}                                      \\
            \frac1{d_{ij}^2}        &\text{(2) N.E.G.} \\
            \exp{-d_{ij}\cdot k}    &\text{(3) exponential decline}
        \end{cases}
\end{align}

This represents a likely more accurate declining effect of productivity spillover,
rather than the special-case group mean where every firm in an enclosed area,
for example, the target's postcode, has an equal influence on the target's TFP.
In addition, it allows for better identification following \cite{Bramoulle2009}.
\todo{explain Bramoulle2009, patterns of }$\Omega_t$
\\*
In essence, the LMM specification only contains 3 non-overlapping groups
from its location tags. The distance-decayed specification
contains as many non-overlapping groups as firms.

\wraptable
    {Static panel using continuous distance-decay weighted TFP average of peers}
    {../../model/output/results_2_dd}
    {}

This distance-decay specification revises down the estimate from peers to a an elasticity of 0.168,
lower than the results given by the LMM specification. This is likely due to \todo{...}
\\
In addition, I introduce controls \todo{...}